\section{Outlook}
\label{main:outlook}

The two applications demonstrate the framework; they do not exhaust it. Both were obtained by the design procedure of Section~\ref{main:chain}, and we see three directions in which it extends.

\paragraph{Unexplored link choices.}
Table~\ref{tab:main-families} is sparse. Every reviewed family builds geometry first and treats the process $\Theta$ as fixed, and the two constructions of this paper fill only two new cells. Sources can lie between a stack of layers and one volume, for instance a few connected pieces chosen to match the part and the machine at once. Traversals of a connected source need not be Hilbert curves; other curves trade locality, smoothness, and turning differently. Construction orders can interleave the links rather than fix one first; a process-aware order, for example, could let thermal or curing constraints in $\Theta$ shape the source split as the posture sheets did here. Each of these choices comes with its own obligations, and the composition rules say in advance what they are.

\paragraph{Lower bounds that combine several links.}
The two bounds of this paper each count components at one link, but obstructions can arise at several at once, as when a part decomposed into $N$ near-convex pieces is printed by a robot whose safe postures form $c$ sheets. Unless the source pieces align with the sheets, we expect the true minimum to be counted by the components of a combined space of source pieces and sheets that a continuous stretch can occupy together, and to exceed both $N-1$ and $c-1$. Stating and proving such bounds, and building sources that attain them, would turn the counting argument from a diagnosis of one link into a design rule for the whole chain.

\paragraph{Uses beyond deposition.}
Any process that covers a region with a prescribed measure along a machine-executed path has the same chain, including machining (material removed), inspection (sensor dwell), and spraying and welding (coating or heat input). Each has its own forced interruptions, such as tool retractions, re-fixturing, and torch restarts, and its own machine with posture sheets.
